\chapter{Methods}

\section{Algorithm Table}

Algorithm \ref{alg:template} demonstrates the thesis-style algorithm table. The algorithm counter is independent from the table counter.
\begin{table}[htbp]
\vspace{0.5em}\centering\wuhao
\refstepcounter{algorithm}
\label{alg:template}
\begin{tabularx}{\linewidth}{@{}lY@{}}
\toprule
\textbf{Algorithm \thealgorithm} & \textbf{Template algorithm table} \\
\midrule
\textbf{Input}
& Initial point $x_0$, stepsize $\gamma$, and maximum iteration number $K$.\\
\vspace{0.2em}
\textbf{Step 1}
& For $k=0,1,\ldots,K-1$, repeat the following computations:\par
\hspace*{2em}1) Compute the gradient $\nabla f(x_k)$;\par
\hspace*{2em}2) Update the variable
$x_{k+1}=\operatorname{prox}_{\gamma g}\!\left(x_k-\gamma\nabla f(x_k)\right)$;\par
\hspace*{2em}3) Check the stopping criterion.\\
\textbf{Output}
& $x_K$.\\
\bottomrule
\end{tabularx}
\end{table}

\section{Theorem-like Environments}

\begin{definition}[Convex function]\label{def:convex}
A function $f:\mathbb{R}^d\to(-\infty,+\infty]$ is convex if, for any $x,y\in\operatorname{dom} f$ and any $t\in[0,1]$,
\begin{equation}\label{eq:convex}
  f(tx+(1-t)y)\le t f(x)+(1-t)f(y).
\end{equation}
\end{definition}

\begin{lemma}\label{lem:basic}
If $f$ is convex, then every local minimizer of $f$ is also a global minimizer.
\end{lemma}

\begin{proof}
The conclusion follows directly from the definition of convexity in Definition \ref{def:convex}.
\end{proof}

\begin{theorem}\label{thm:template}
Under suitable assumptions, the sequence generated by Algorithm \ref{alg:template} is well defined.
\end{theorem}

\begin{proof}
This is a template proof. Replace it with the actual proof in the thesis.
\end{proof}

\begin{remark}
The theorem, lemma, definition, assumption, corollary, proposition, example, and remark environments are numbered independently by chapter.
\end{remark}
